\documentclass[conference]{IEEEtran}

\usepackage[utf8]{inputenc}
\usepackage{graphicx}
\usepackage{hyperref}
\usepackage{booktabs}
\usepackage{cite}

\title{MedGov-AI: An MCP-Based Agentic Orchestrator for Clinical Medical Image Analysis and Structured Report Generation}

% too current text? maybe we should do more separation or starting bolding and puting in italic some things


\begin{document}
\maketitle

\begin{abstract}
\end{abstract}

\begin{IEEEkeywords}
Agentic AI, Medical Imaging, MCP, MONAI, Clinical Workflow Automation
\end{IEEEkeywords}

\section{Background}

\subsection{AI in Medical Imaging}
% talk about AI in medical imaging
% real case scenarios and uses
% risks
% pros and cons
Before starting with all the technologies, methods and results, we need to understand why this work was started in the first place, what was the necessity that brought this possible solution, does it work? Is it something still in development and being explored? Are new technologies really necessary for this line of work? These and more will all be questions that we hope to answer with our experience in this project.
% get real metrics for this
Healthcare systems worldwide face increasing demand with limited resources. The volume of medical imaging studies continues to grow while the number of specialists remains constrained, creating diagnostic backlogs that directly affect patient outcomes. Automating routine analysis tasks is therefore not merely a convenience but a structural necessity. Greater automation brings greater consistency, reduces the workload on healthcare professionals, accelerates diagnosis, and allows clinicians to dedicate more time to complex cases and direct patient care.
Since the project focuses on the exploration of AI and MCP in order to provide assistance in medical fields we need to ask these 4 questions:
% this part will be filled more as we read some papers in order to answer, and I dont know if the answer should be right after on ir another sections

- What AI tools are currently being used in medical imaging? (Segmentation models, detection, report generation, NLP on clinical notes)
% research had to be done
Several AI tools are currently being used in medical imaging, with the notable detail that they typically operate in isolation from one another, with no single standard way of connecting them. Examples include segmentation and detection systems such as chest X-ray analysis for detecting pneumonia, nodules and fractures used in triage, CT scan analysis for lung nodule detection with FDA-cleared tools such as Viz.ai and Aidoc, mammography screening with AI-assisted breast cancer detection tools such as Transpara and iCAD, and diabetic retinopathy screening from fundus images with systems such as IDx-DR.
There are also tools for radiology reporting, such as automated preliminary reads that flag findings for radiologist review, structured report generation, and critical finding alerts. In pathology, cell counting and tissue classification models are in active clinical use. More general capabilities include worklist prioritization, where AI decides which scan a radiologist should read first based on predicted urgency.

- What are they actually being used for? (Workflow automation, decision support, quality control, radiology, pathology, etc.)
% search for hospitals admitting to the use of AI - verify all examples below before final submission
These tools are being used across several areas of clinical practice, with adoption documented in leading institutions worldwide. In radiology, AI assists with worklist prioritization and critical finding alerts, where the system detects life-threatening conditions such as stroke, pulmonary embolism or intracranial hemorrhage and immediately notifies the clinical team. Mass General Brigham (United States) has deployed AI-driven radiology worklist prioritization to reduce time-to-read for urgent studies. The NHS (United Kingdom) deployed Annalise.ai across more than 40 hospital trusts for chest X-ray triage, allowing radiologists to focus on flagged abnormal cases first. In screening programs, AI acts as a first filter, flagging high-risk cases for specialist review and allowing programs to scale without proportionally increasing the number of specialists required. Google Health conducted a published clinical trial of an AI system for diabetic retinopathy screening in Thailand, demonstrating that AI-assisted grading matched or exceeded specialist performance in a real clinical setting. Mayo Clinic (United States) has integrated AI into ECG interpretation and cardiac imaging workflows, using it as a decision support layer for cardiologists. In pathology, whole slide image analysis is used for tumour grading and cell counting, reducing the time needed to evaluate complex specimens. More broadly, AI is being used for documentation support, generating structured preliminary reports from imaging findings that the clinician then reviews and validates.
% --- references to add ---
% Mass General Brigham + Annalise.ai deployment:
% https://www.massgeneralbrigham.org/en/about/newsroom/press-releases/annalise-ai-imaging-tools
% https://www.healthcareitnews.com/news/mass-general-brigham-and-future-ai-radiology
% NHS Annalise.ai across 40+ trusts:
% https://harrison.ai/news/transformational-ai-diagnostic-tool-made-available-to-radiologists-in-over-40-nhs-trusts/
% https://www.leedsth.nhs.uk/news/new-ai-chest-x-ray-software-to-speed-up-diagnosis-and-improve-patient-care-at-leeds/
% Google Health Thailand trial (The Lancet Digital Health - add as citation):
% https://www.thelancet.com/journals/landig/article/PIIS2589-7500(22)00017-6/fulltext
% https://pubmed.ncbi.nlm.nih.gov/35272972/
% NOTE: the Thailand trial also found real-world challenges worth mentioning in risks/limitations:
% image quality requirements that staff could not always meet, inconsistent workflows between
% clinics, and connectivity issues - good example that real-world deployment is harder than lab results
% Mayo Clinic AI ECG:
% https://www.mayoclinic.org/departments-centers/ai-cardiology/overview/ovc-20486648
% https://newsnetwork.mayoclinic.org/discussion/spotlight-on-early-detection-of-3-heart-diseases-using-ecg-ai/
% Viz.ai, Aidoc, Transpara, iCAD, IDx-DR - to add references after
- What are the benefits and why could this be a solution? (faster diagnosis, reduced workload, consistency)
The benefits of AI in medical fields are directly tied to the limitations of the current manual process. Faster diagnosis is perhaps the most immediate gain, as AI systems can process and flag findings in seconds, reducing the time between scan acquisition and clinical decision. Consistency is another key advantage, although not fully achieved in today's world, particularly when referring to general large language models and chatbots. However, as these systems become more widely used in clinical settings and more medically specific, they will only improve. Human specialists on the other hand are subject to many influencing factors such as fatigue, experience, and personal circumstances, whereas AI models apply the same criteria uniformly across every case. The reduction in workload is also significant, as automating routine reads and preliminary reports frees specialists to focus on cases that genuinely require expert judgement. Taken together, these benefits suggest that AI is not a replacement for clinicians but a force multiplier.

- What are the risks and limitations? (bias, hallucinations, black box decisions, regulatory concerns, no existing standards for agentic systems in healthcare)
The risks and limitations of AI in medical imaging are as significant as its potential benefits. One of the most pressing concerns is bias, as AI models are only as representative as the data they were trained on, meaning that systems trained predominantly on data from certain demographics or imaging equipment may underperform on patients from underrepresented groups or institutions using different hardware and protocols. Hallucination is another critical risk, particularly in large language models used for report generation, where the system may produce confident but clinically incorrect statements. The black-box nature of many deep learning models raises accountability concerns, since clinicians and regulators often cannot explain why a model reached a particular conclusion, which complicates both clinical trust and legal responsibility. From a regulatory standpoint, the approval and monitoring of AI tools in healthcare remains inconsistent across countries and institutions, with no unified framework governing how these systems should be validated, updated, or audited in production. Finally, and most relevant to this work, there are currently no established standards for how agentic AI systems should operate in clinical environments, including how they should communicate with existing tools, when they should escalate to a human, and how their decisions should be logged for audit.


Despite the progress of individual AI tools, no standard integration layer exists to orchestrate them together within a clinical workflow. Each institution must develop its own connectors and its own standards, which leads to fragmentation, lack of interoperability, slow adoption, and deployments that are expensive and difficult to maintain. The Model Context Protocol (MCP) is a promising solution to this problem, providing a standardized way for agents to discover and call tools without bespoke integration. This work explores the application of MCP as the integration layer for an agentic AI system that orchestrates multiple clinical imaging tools in a real-world workflow.

\subsection{Agentic AI in Healthcare}

% define agentic AI: goal-directed, tool-using, multi-step reasoning
% why an agent fits clinical imaging workflows (sequential by nature)
% exploration of agentic systems that we tried and conclusions
% mention maybe also safety concerns, more adaptability, the customization of the system and determined specific clinical workflows, and the fact that we wanted to explore the design space of agentic systems in healthcare more broadly rather than being constrained by the design choices of existing frameworks.

Agentic AI refers to systems that can autonomously set goals, plan multi-step actions, and use tools to achieve those goals. Unlike traditional AI models that operate in a single step or require human input at every stage, agentic systems can reason through complex tasks, make decisions about which tools to use, and adapt their behavior based on feedback. This makes them particularly well-suited for clinical imaging workflows, which are inherently sequential and often require the integration of multiple tools and data sources. For example, a typical workflow for analyzing a CT scan might involve first running a segmentation model to identify regions of interest, then using a classification model to characterize those regions, and finally generating a structured report summarizing the findings. An agentic AI system can orchestrate this entire process autonomously, calling each tool in the correct sequence and synthesizing the results into a coherent output.

Three core properties characterize agentic systems: goal persistence (the system does not stop until the required deliverables are met), tool use (it can act on external systems, not just generate text), and iterative reasoning (each result informs the next decision). We evaluated several existing agentic frameworks, including LangChain, CrewAI, and AutoGen, but opted to build a custom agentic system. The primary reason was the need for tight integration with MCP: existing frameworks do not natively support MCP as a tool protocol, which would have required additional adapter layers and reintroduced the fragmentation problem we were trying to solve. Building a custom loop also allowed us to enforce specific behaviors required in a clinical context, such as a completion contract under which the agent cannot declare success while required deliverables remain unmet, and a debug mode in which the user must approve each tool call before execution.

Contrasting with simpler alternatives clarifies why this design was necessary. A plain chatbot produces one response per prompt and cannot maintain state or call tools. A hardcoded pipeline can call tools but relies on the developer to anticipate every possible execution path, making it brittle in the face of real-world variability. A standalone ML model provides predictions but has no language interface and cannot adapt its behavior based on context. We found that only a true agentic system could handle the complexity and variability of real clinical workflows, which often require conditional logic, error handling, and dynamic tool selection.

Our own system began as precisely this kind of hardcoded pipeline: in the initial version, the LLM was never consulted, and all sequencing decisions were made by explicit conditional logic in the orchestrator code. This approach proved limiting, as it could not adapt to new information without significant reprogramming, which motivated the transition to a fully agentic design.

This transition is not without its challenges. Reliable agentic behavior requires careful prompt engineering: the agent must be guided to select tools in the correct order, avoid premature declarations of success, and handle ambiguous or incomplete results gracefully. These challenges are discussed in detail in later sections. However, the fundamental advantage over any hardcoded alternative is scalability: supporting a new clinical domain does not require rewriting workflow logic but only adding a new MCP server that exposes the relevant tools, or a skill file that encodes the domain-specific protocol, and the agent discovers and incorporates both at runtime without code changes. This also opens a much wider architectural design space. Rather than a single fixed agent, one could configure specialist agents for distinct clinical tasks, use prompt-based behavioral switches to change the agent's operating mode, or compose multiple MCP servers covering entirely different domains. None of these extensions require touching the orchestration code, which is simply not true of any hardcoded pipeline.

% Note: LLM backbone comparison placed here; consider moving to Methods or Results if the teacher prefers

% is it worth mentioning the hardware it was tested on? and enumerating all the models for example?
When testing the system we explored two LLM backend approaches: a locally-hosted model and a cloud API. The local model was Ollama's deepseek-r1:7b, with llama3.1:8b also evaluated, both small enough to run on consumer hardware and allowing for complete data privacy since all processing happens locally. The API-based model was Gemini 2.0 Flash, with later versions used as token limits required, a significantly larger and more capable model hosted in the cloud. The local approach offers greater control and privacy but proved to struggle with the demands of multi-step clinical workflows. We observed consistent failures in tool calling: the model frequently failed to select the correct tool for the given context and could not reliably chain multiple calls to reach a goal. Addressing this consumed considerable development time across prompt engineering, workflow simplification, logging, and targeted adjustments. None of these fully resolved the problem, as the root cause was model size rather than configuration. Switching to the API-hosted model resolved all of these issues immediately. The API approach introduces trade-offs of its own, namely latency and potential data security concerns in privacy-sensitive deployments, but for the complex reasoning required by clinical imaging workflows it proved to be the only reliably viable option at the model sizes currently available for local inference. The system architecture supports both backends, allowing deployment to be tailored to the constraints of a given clinical environment.

\subsection{MCP as a Standardized Tool Integration Layer}

The conventional approach to agent-tool integration is the bespoke REST adapter: for each tool the agent must call, a custom connector is written that handles authentication, request serialization, response parsing, and error mapping. This approach does not scale. Each new tool adds a new adapter; changes to a tool API break the corresponding adapter; and the agent itself must encode assumptions about every tool it might call.

The Model Context Protocol (MCP) addresses this directly. MCP defines a client-server protocol in which each tool server declares its available tools at startup, specifying names, input schemas, and descriptions. The agent discovers tools dynamically at runtime and invokes them by name without requiring any prior knowledge of the server's internal implementation. This decouples the agent from individual tool APIs: adding a new tool server requires no changes to the agent, only a new MCP-compliant server that exposes the relevant capabilities.

Singh et al. (2025) survey MCP's architecture in detail, positioning it as the solution to fragmented API integrations for agentic AI systems~\cite{mcp-aditi}. The survey also identifies MCP sampling, a mechanism for LLM-to-LLM communication within the MCP protocol, as a path toward sub-agent specialization, in which a coordinating agent can delegate to specialist agents exposed as MCP servers. Singh et al. note that healthcare is a domain where these properties are especially relevant, but the survey discusses clinical applications only conceptually; no implemented system is presented.

\subsection{Skills as Clinical Capability Packages}

Tool access via MCP determines what an agent can do, it does not determine what the agent should do in a given clinical domain. A segmentation server, a terminology server, and a report template server can all be exposed as MCP tools, but the correct sequencing of those tools for a CT lung evaluation versus a pathology workflow requires domain knowledge that is separate from both the agent's general reasoning capability and the tool's own logic.

Skills address this layer. Xu and Yan (2026) formalize agent skills as composable, on-demand capability packages that provide domain-specific behavioral guidance to an LLM without retraining it~\cite{agentskillslargelanguage}. The key architectural insight is that skills and MCP are complementary rather than competing: skills define what to do in a given domain, while MCP exposes the tools with which to do it. Xu and Yan do not address healthcare-specific skill design, the governance question of who maintains skill files when clinical protocols change, or how skills compose with MCP in a working system, gaps that this work directly addresses.

\subsection{Positioning of This Work}
% no prior work combines: DICOM parsing + segmentation + reports under one agent,
% MCP as the actual interoperability layer, and skill files for protocol encoding
% reference the MCP-FHIR paper which found the exact same fragmentation problem

\section{Methods}

\subsection{System Architecture Overview}

\begin{figure}[t]
  \centering
  \includegraphics[width=\columnwidth]{images/architecture.png}
  \caption{MedGov-AI system architecture.}
  \label{fig:arch}
\end{figure}


MedGov-AI is organized into three layers: a React frontend, a FastAPI
orchestrator, and a set of MCP tool servers. Each layer has a clearly
defined responsibility, and the boundaries between them are designed so
that clinical domain logic lives neither in the interface nor in the
transport layer but entirely in the orchestrator, where it can be
reasoned about, logged and audited(to implement ******).

\subsubsection*{Frontend}
The client layer is a React~18 application that provides the clinical
user interface, communicating with the orchestrator exclusively via a
REST and SSE API\@. SSE enables real-time streaming of task progress
and agent responses to the user without polling.

\subsubsection*{Orchestrator}
The orchestrator is a FastAPI application and is the system's single
point of control: it exposes the REST and SSE API consumed by the
frontend, owns all session and task state, manages connections to every
MCP server, and hosts the three internal services that drive the
system's behaviour: the \texttt{AgenticAgent}, the
\texttt{task\_runner}, and the \texttt{watcher\_service}.

The \texttt{AgenticAgent} is the sole decision-maker in the system. A
deliberate choice was made to use a single agent rather than a
hierarchy of sub-agents, for two reasons: simplicity and auditability.
A single agent is easier to reason about, debug, and extend, and every
tool call, its inputs, and its result are recorded in a single session
trace with no hidden delegation. The agent runs an iterative LLM loop
using either the Gemini API or a locally-hosted Ollama model as its
reasoning backend, and executes multi-turn function calling to invoke
MCP tools in sequence. It maintains a session context across turns so
that information produced by one tool call is available to the next
without requiring the user to restate it.

The \texttt{task\_runner} manages long-running operations that cannot
complete within a synchronous HTTP request. Operations such as MONAI
inference or Cellpose segmentation are submitted to a background thread
pool and tracked through a state machine with four states: queued,
running, completed, and failed. A semaphore limits concurrent GPU
inference to a single job at a time, preventing out-of-memory crashes
on constrained hardware. The task runner also captures stderr output
from inference processes and, on failure, passes the log to the LLM
for a plain-language error explanation before surfacing it to the user.
Task state is persisted in SQLite so that progress survives backend
restarts, and the frontend is notified of state transitions via SSE\@.

The \texttt{watcher\_service} monitors registered filesystem
directories using periodic polling and triggers the agent automatically
when a new file is deposited. This enables event-driven clinical
workflows in which files arriving from a PACS system or a scanner
initiate the full analysis pipeline without any user action. No changes
to upstream systems are required; the watcher integrates at the
filesystem level and streams its activity log to the frontend via a
dedicated SSE channel.

All persistent state is stored in a single SQLite database with WAL
journaling. The schema covers users, sessions, conversation history,
uploaded files, background task status, watched directories, and
per-user tool settings.

\subsubsection*{LLM Backends}
The orchestrator supports two LLM backends, selectable at deployment
time without modifying any agent logic. The Gemini API provides a
large cloud-hosted model capable of reliable multi-step function
calling in complex clinical workflows, at the cost of external data
transmission and API dependency. Ollama provides a local inference
alternative using smaller open-weight models running entirely on the
deployment host with no external data transmission, which is relevant
in privacy-sensitive clinical environments, but proved insufficient for
reliable tool calling at the 7--8B parameter scale currently available
for local inference. Both backends implement the same internal
interface.

\subsubsection*{MCP Server Layer}
The tool layer consists of seven MCP servers, each running as an
independent process in its own isolated Python virtual environment.
Isolation at the process level is a deliberate design choice: it
prevents dependency conflicts between servers with incompatible
requirements (most critically between the GPU-enabled PyTorch
installation required by \texttt{mcp-monai} and the CPU-only
installation used elsewhere, which cannot coexist in a single
environment). Tool names are prefixed with the server name
(e.g., \texttt{monai.run\_inference}, \texttt{radlex.generate\_report}),
providing unambiguous namespacing across servers. Six servers
communicate over stdio transport; the FHIR server uses HTTP transport
to connect to an external endpoint. The available servers and their
clinical functions are listed in Table~\ref{tab:mcp-servers}.

\begin{table}[t]
  \caption{MCP Servers and Clinical Functions}
  \label{tab:mcp-servers}
  \centering
  \begin{tabular}{lll}
    \toprule
    Server & Transport & Clinical Function \\
    \midrule
    mcp-monai    & stdio & Medical image segmentation and detection \\
    mcp-radlex   & stdio & RadLex terminology and report templates \\
    mcp-utils    & stdio & DICOM parsing and filesystem operations \\
    mcp-skills   & stdio & Skill workflow loading and execution \\
    mcp-ipath    & stdio & iPath telepathology integration \\
    mcp-cellpose & stdio & 2D/3D cell segmentation (Cellpose) \\
    mcp-fhir     & HTTP  & FHIR patient data retrieval \\
    \bottomrule
  \end{tabular}
\end{table}

This architecture means that extending the system with a new clinical
capability requires only writing a new MCP-compliant server and
registering it in the configuration file. The agent, the orchestrator,
and every existing server remain unchanged. This is the property that makes MCP a viable interoperability layer rather than yet another bespoke integration approach.


\subsection{MCP Server Layer}

Each MCP server encapsulates a distinct clinical capability and exposes
it as a set of named, schema-defined tools. The agent has no hardcoded
knowledge of any server's internals. It discovers available tools at
startup and selects among them based on context.

\textbf{Mcp-monai} is the medical image inference server, built on the
MONAI framework~\cite{monai} with PyTorch as the underlying deep
learning backend. \texttt{analyze\_image} detects the imaging modality
and body part automatically from pixel statistics and file metadata,
and returns a ranked list of recommended models. \texttt{list\_models}
filters the model catalogue by modality, body part, or category.
\texttt{download\_model} retrieves a MONAI bundle from HuggingFace Hub
and caches it locally. \texttt{run\_inference} executes a segmentation
or detection model using a sliding window strategy to handle the
variable dimensions of 3D medical volumes, and submits the job to the
background task queue to avoid blocking the session.
\texttt{get\_monai\_info} returns the current MONAI and PyTorch version
and GPU availability. Bundled models cover spleen segmentation on CT,
pancreas and tumour segmentation on CT, lung nodule detection on CT,
brain tumour segmentation on multi-channel MRI, whole-body CT
segmentation across 104 anatomical structures, and kidney segmentation
using the UNEsT architecture.

\textbf{Mcp-radlex} provides structured radiology reporting through
integration with the RadLex terminology standard.
\texttt{list\_subspecialties} returns available radiological
subspecialties. \texttt{find\_templates} searches the template
catalogue by subspecialty or keyword. \texttt{get\_template\_schema}
fetches the full field structure of a given template, including
required and optional fields. \texttt{generate\_report} populates a
template with agent-provided findings and returns a structured,
validated report. This server closes the loop between inference output
and clinical documentation: the agent receives segmentation results,
interprets them, and produces a report conforming to a standard
radiological schema rather than free text.

\textbf{Mcp-utils} provides general-purpose filesystem and DICOM
operations that underpin every clinical workflow. \texttt{parse\_dicom}
extracts the full metadata of a single DICOM file (modality, body
part, patient identifiers, acquisition parameters, and pixel data
statistics) while \texttt{parse\_dicom\_directory} applies the same
extraction across an entire directory. On the filesystem side, the
server exposes tools for reading, writing, moving, copying, and
deleting files, listing directory contents, and searching by pattern.

\textbf{Mcp-cellpose} provides cell segmentation for pathology
workflows using the Cellpose universal model (cpsam, v4). It exposes
segmentation tools for both 2D and 3D inputs, handling the different
dimensionalities encountered in histology and volumetric imaging.
\texttt{save\_overlay} composites the predicted cell mask onto the
original image and writes the result to disk. GPU availability is
detected at startup with automatic CPU fallback.

\textbf{Mcp-ipath} integrates with the iPath telepathology platform,
which aggregates whole slide images (WSI) and DICOM series via
Dicoogle and DICOMweb. \texttt{fetch\_thumbnail} retrieves a scaled
preview of a WSI by UID\@. \texttt{get\_slide\_dimensions} returns full
pixel dimensions. \texttt{scale\_roi\_to\_slide} converts a
region-of-interest across resolution levels. \texttt{fetch\_roi}
retrieves a specific region at a given resolution without downloading
the full WSI\@. \texttt{fetch\_dicom\_instance} and
\texttt{get\_series\_instances} provide access to individual DICOM
frames and series.

\textbf{Mcp-skills} exposes the skill system to the agent as callable
tools. \texttt{read\_skill\_file} loads a skill's entry point and
returns its full content. \texttt{read\_references} retrieves a
reference document stored alongside a skill.
\texttt{read\_asset} returns an asset associated with a skill, such as
a report template. \texttt{execute\_script} runs a shell command
defined within a skill and returns its output. The skills system is
described in detail in Section~\ref{sec:skills}.


\subsection{Skills System}
\label{sec:skills}

MCP defines what tools are available to the agent; it does not define
how those tools should be used in a given clinical domain. The correct
sequence of operations for a CT lung workflow (parse DICOM, detect
modality, select an appropriate segmentation model, run inference,
retrieve results, fill a RadLex report template) is not derivable
from tool schemas alone. It requires domain knowledge about clinical
protocols, sequencing constraints, and the interpretation of
intermediate results. The skills system encodes this knowledge in a
form the agent can load on demand without any retraining.

This knowledge could in principle be embedded directly in the system
prompt, but doing so has two drawbacks. First, it increases token
consumption on every request, including those that do not require that
domain. Second, a large system prompt increases the risk of the LLM
overlooking or misapplying specific instructions, a well-documented
failure mode in prompt engineering~\cite{}. Skills address both
problems by keeping the system prompt minimal and loading
domain-specific protocol only when the agent determines it is needed
for the current task.

\subsubsection*{Structure}

Each skill is a directory containing a \texttt{SKILL.md} file as its
entry point, written in YAML-frontmatter Markdown. The frontmatter
defines the skill's name and a natural language description that is
injected into the agent's system prompt at startup. When the agent
encounters a task that matches a skill's description, it calls
\texttt{skills.read\_skill\_file} to load the full content before
taking any action in that domain. The skill file then acts as the
agent's operating instructions for that workflow.

The skill directory can contain additional resources alongside
\texttt{SKILL.md}: reference documents that detail specific
sub-workflows, assets such as report templates or annotated examples,
and executable scripts that can be invoked via
\texttt{skills.execute\_script}. This structure allows a single skill
to encode multiple related workflows, with \texttt{SKILL.md} acting as
a router that directs the agent to the appropriate reference file based
on the user's request. This filesystem-based architecture enables
progressive disclosure: the agent loads only the router at first, then
fetches deeper protocol detail on demand, keeping the context window
focused and avoiding unnecessary token consumption on workflows that
are not relevant to the current task.

\subsubsection*{Implemented Skills}

Two skills are currently deployed in the system. The
\textbf{medical-imaging} skill covers any task involving medical image
data. Its \texttt{SKILL.md} functions as a workflow router: it presents
two available workflows (DICOM analysis with report generation, and
whole slide image tumour detection) and instructs the agent to load
the corresponding reference document before proceeding. Each reference
document contains the complete step-by-step protocol for that workflow,
including which tools to call, in what order, how to handle errors, and
what the expected outputs are. This separation between the router and
the protocol documents means that new imaging workflows can be added to
the skill without modifying the router or any code.

The \textbf{clinical-reports} skill covers structured documentation
across a range of report types, including radiology and pathology
diagnostic reports, case reports following CARE guidelines, clinical
trial reports following ICH-E3 and related regulatory standards, and
patient documentation formats such as SOAP notes, H\&P documents, and
discharge summaries. It encodes compliance requirements (HIPAA, FDA,
ICH-GCP) and provides templates and validation guidance that the agent
applies when generating documentation. Unlike the medical-imaging
skill, which was developed specifically for this work, the
clinical-reports skill is a publicly available skill integrated without
modification, demonstrating that third-party skills can be adopted
into the system without any code changes.

\subsubsection*{Governance}

A key property of this design is that skill files are plain text and
can be updated by domain experts (clinicians, radiologists,
pathologists) without touching the agent, the MCP servers, or any
other part of the codebase. When a clinical protocol changes, the
relevant reference document is updated and the agent adopts the new
behaviour immediately on the next session, with no software deployment
required. This operationalises the skills-as-updatable-protocol
concept described by Xu and Yan~\cite{agentskillslargelanguage} in a
concrete clinical setting, and directly addresses the governance gap
they identified: the question of who maintains skill files and how they
are kept current with clinical practice.


\subsection{Agentic Execution Loop}
% goal in, tools called iteratively, result evaluated each step
% completion contract: agent cannot declare success while required deliverables are missing
% two modes: autonomous (no interruption) and debug (confirmation before each tool call)
% LLM backbone configurable: Gemini (cloud) or Ollama (local, privacy-preserving)

\subsection{Background Task Queue}
% long-running operations (MONAI inference, Cellpose) queued to thread pool
% state machine: queued, running, completed/failed, persisted in SQLite
% SSE notifications to frontend on completion
% semaphore: only one inference runs at a time to prevent memory crashes

\subsection{Workspace Monitoring}
% named monitored directories registered in the system
% file deposited triggers agent automatically with a predefined prompt
% event-driven clinical workflows: files from PACS or scanner initiate pipeline
% no change to upstream systems required

\subsection{Evaluation Setup}
% the 3 experimental phases defined in the March 2026 meeting with the teacher
% hardware: NVIDIA GeForce RTX 3050 6GB, MONAI 1.5.1, PyTorch 2.5.1
% data: public DICOM datasets (Cancer Imaging Archive, Medical Segmentation Decathlon)
% LLM backends tested: Gemini 2.0/2.5 Flash (API), Ollama deepseek-r1:7b and llama3.1:8b

\section{Results}

\subsection{Phase 1 - Model Experimentation}
% Gemini vs Ollama for tool calling in clinical workflows
% local versus API, limitations, pros and cons

\subsection{Phase 2 - MCP Tools and Skills}
% tool discovery at startup, dynamic registration
% skills experimentation: advantages, disadvantages, findings

\subsection{Phase 3 - Use Case Comparison}
% use cases go here
% autonomous workflow vs manual/guided workflow

\section{Discussion}

\section{Conclusions}

\bibliographystyle{IEEEtran}
\bibliography{references}

\end{document}
